\subsection{Limitaciones interpretativas}

\begin{itemize}[leftmargin=1.27cm]
  \item El CNPV mide condiciones de 2018; no representa por sí solo el contexto socioeconómico de 2026.
  \item La unidad territorial y el cambio de límites entre 2018 y 2026 generan una incompatibilidad potencial. Los códigos 27493 y 94663 requieren crosswalk verificable antes de comparar contratación.
  \item Las veinte ANM de 2018 se identifican y se describen aparte. El ámbito descriptivo principal incluye municipios/distritos y San Andrés; su inclusión insular en el futuro PCA requiere decisión explícita.
  \item Hogares con menores y con personas mayores son proxies de carga potencial de cuidado, no mediciones de dependencia funcional, económica o de necesidad de cuidado.
  \item 5PM agrupa situaciones heterogéneas y usa una ponderación normativa 2:1. El indicador no es una tasa oficial de desempleo ni describe actividad de un hogar.
  \item El estrato es una clasificación del inmueble para servicios, no una medida directa de ingreso. La inclusión de \enquote{No tiene estrato} como condición de vulnerabilidad es una decisión del equipo, cuya sensibilidad al denominador es alta.
  \item 17PM usa población 5+ y no controla composición de edad. Se excluye del primer PCA y se reserva para sensibilidad exploratoria.
  \item La matriz z equaliza varianzas, pero no corrige asimetría, observaciones influyentes, error conceptual o solapamiento entre indicadores.
  \item No se han actualizado KMO, Bartlett, MSA, análisis paralelo ni estabilidad para las ocho candidatas y el universo municipal final.
  \item SECOP II no se ha integrado. La comparación 2018--2026 será temporalmente asimétrica y ecológica; los resultados describirán asociación, no causalidad.
\end{itemize}

\subsection{Decisiones adoptadas y asuntos abiertos}

\begin{table}[htbp]
\centering\small
\caption{Estado de decisiones metodológicas}\label{tab:decisiones_pendientes}
\begin{adjustbox}{max width=\textwidth}
\begin{tabular}{@{}p{0.32\textwidth}p{0.28\textwidth}p{0.32\textwidth}@{}}
\toprule
\textbf{Decisión / asunto} & \textbf{Estado} & \textbf{Siguiente acción}\\
\midrule
IVS con indicadores DANE; SECOP como análisis posterior & Acordado & Mantener separados los bloques de vulnerabilidad y contratación. \\
Base censal 2018; SECOP II prioriza 2026 & Acordado & Documentar desfase temporal y cobertura. \\
ANM por fuera del IVS municipal principal & Lista histórica verificada; base separada & Cerrar tratamiento de San Andrés para el PCA. \\
27493 hacia 27615 para armonización & Decisión de enlace del equipo & Verificar límites/fuente oficial; no alterar 2018. \\
94663 hacia 94343 para armonización & Decisión de enlace del equipo & Agregar componentes de 2018 a geografía común antes de tasas, si el diseño de enlace lo requiere. \\
Cobertura SECOP 100\% del universo acordado & Criterio del equipo & Fijar denominador después de resolver universo y crosswalk. \\
5PM peso 2 para incapacidad & Aprobado como ponderación experimental & Reportar sensibilidad 1:1 y 3:1; no presentarlo como peso empírico. \\
17PM fuera del primer PCA & Acordado & Comparar luego solución de ocho variables con variante educativa, si se valida su universo. \\
Basuras fuera de las ocho & Confirmado el 18 de octubre & Conservar únicamente como sensibilidad exploratoria. \\
Solapamiento con IPM & Criterio acordado; matriz conceptual documentada & Contrastar después soluciones con retiro de candidatas cercanas. \\
PCA / IVS / validación externa / SECOP & Pendientes & Ejecutar en etapas, con diagnósticos y trazabilidad versionados. \\
\bottomrule
\end{tabular}
\end{adjustbox}
\end{table}

\subsection{Próximas etapas}

\begin{enumerate}[leftmargin=1.27cm]
  \item Cerrar la inclusión insular en el universo del PCA y verificar el crosswalk 2018--2026; la clasificación de veinte ANM ya está documentada.
  \item Seleccionar el ámbito definitivo desde las matrices de ocho variables y volver a estandarizarlo; preservar el escenario separado para ANM.
  \item Calcular KMO/MSA y Bartlett cuando se autorice la etapa multivariada. Los controles de fórmula, descriptivos, correlaciones y cotejo de extremos del ámbito municipal e insular ya se documentan en esta actualización.
  \item Ejecutar PCA exploratorio, comparar criterios de retención e inspeccionar estabilidad de cargas y puntajes bajo sensibilidades previamente definidas.
  \item Si la solución es interpretable y estable, formalizar la regla de construcción del IVS; posteriormente contrastarlo con IPM como criterio externo.
  \item Obtener SECOP II 2026, documentar cobertura municipal y armonizar límites antes de analizar asociación, correspondencia y brechas.
\end{enumerate}

La secuencia es consistente con guías de construcción de indicadores compuestos: las decisiones de normalización, ponderación, agregación y análisis de incertidumbre deben ser transparentes y evaluarse antes de interpretar un ranking \citep{oecd2008composite}. El IVS no se considerará terminado solo por obtener una primera componente principal.
